%%%%%%%%%%%%%%%%%%%%%%%%%%%%%%%%%%%%%%%%%%%%%%%%%%%%%%%%%%%%%%%%%%%%%%
% Lớp 11 | Quan hệ vuông góc trong không gian | Đường thẳng vuông góc với mặt phẳng
% Dạng 1H8?2-1: Câu hỏi lý thuyết
% Số câu: 39  (sinh tự động bởi tach_cau.py; từ nay thầy sửa trực tiếp file này)
%%%%%%%%%%%%%%%%%%%%%%%%%%%%%%%%%%%%%%%%%%%%%%%%%%%%%%%%%%%%%%%%%%%%%%

%%%%%%%%%%%%%%%% TRẮC NGHIỆM 4 PHƯƠNG ÁN %%%%%%%%%%%%%%%%

% nguồn: Ok_Toan11_TracNghiem_Lan2
\begin{ex}%[1H8N2-1]%[ma:272cc302]
	Khẳng định nào  sau đây là \textbf{sai}?
	 \choice
		{Nếu đường thẳng $d\perp (\alpha)$ thì $d$ sẽ vuông góc với  mọi đường thẳng nằm trong $(\alpha)$}
		{\True Nếu đường thẳng $d$ vuông góc với hai  đường thẳng nằm trong $(\alpha)$ thì $d\perp (\alpha)$}
		{Nếu đường thẳng $d$ vuông góc với hai đường thẳng cắt  nhau nằm trong $(\alpha)$ thì $d$ vuông góc với mọi đường thẳng nằm trong $(\alpha)$}
		{Nếu đường thẳng $d\perp (\alpha)$ và $a\parallel (\alpha)$ thì $d\perp  a$}
	\loigiai{ 
		Điều kiện cần và đủ để $d\perp (\alpha)$ là $d$ vuông góc với hai đường thẳng cắt nhau nằm trong $(\alpha)$.\\ 
		Vậy khẳng định sai là \lq\lq Nếu đường thẳng $d$ vuông góc với hai đường thẳng nằm  trong $(\alpha)$ thì $d\perp (\alpha)$\rq\rq.
	 }
\end{ex}

% nguồn: Ok_Toan11_TracNghiem_Lan2
\begin{ex}%[1H8N2-1]%[ma:fc6fd88a]
	Chọn mệnh đề đúng  trong các mệnh đề sau.
	 \choice
		{Góc giữa  hai mặt phẳng bằng góc giữa hai đường thẳng tùy ý nằm trong mỗi mặt phẳng}
		{\True Góc giữa hai mặt phẳng bằng góc giữa hai đường thẳng  lần lượt vuông góc với hai mặt phẳng đó}
		{Góc  giữa hai mặt phẳng luôn là góc nhọn}
		{Góc giữa hai mặt phẳng bằng góc giữa hai vec tơ chỉ phương của hai đường thẳng lần lượt vuông góc với hai  mặt phẳng đó}
	 \loigiai{
		Mệnh đề đúng là \lq\lq Góc giữa hai mặt phẳng bằng góc giữa  hai đường thẳng lần lượt vuông góc với hai mặt phẳng đó\rq\rq.
	 }
\end{ex}

% nguồn: Ok_Toan11_TracNghiem_Lan2
\begin{ex}%[1H8N2-1]%[ma:3be9a1a1]
	Cho mặt phẳng $(P)$  và hai đường thẳng phân biệt $a$ và $b$. Biết rằng $a\parallel (P)$. Hỏi mệnh đề nào dưới đây đúng?
	\choice 
		{Nếu $b\parallel (P)$ thì $b\parallel a$} 
		{\True  Nếu $b\perp (P)$ thì $b\perp a$}
		 {Nếu $b\parallel a$ thì $b\parallel (P)$}
		{Nếu $b\perp a$ thì  $b\perp (P)$}
	 \loigiai{
		Ta có nếu $\heva{& a\parallel (P)\\& b\perp (P)}$ thì  $a\perp b$.\\
		Vậy mệnh đề đúng là \lq\lq Nếu $b\perp (P)$ thì  $b\perp a$\rq\rq.
	 }
\end{ex}

% nguồn: Ok_Toan11_TracNghiem_Lan2
\begin{ex}%[1H8N2-1]%[ma:29b213b8]
	Cho hai đường thẳng $a$, $b$ và mặt phẳng $(P)$. Trong các mệnh đề sau mệnh đề  nào \textbf{sai}?
	 \choice
		{Nếu $a\parallel  (P)$ và $b\perp (P)$ thì $a\perp b$}
		{Nếu  $a\subset (P)$ và $b\perp (P)$ thì $a\perp b$}
		{Nếu  $a\perp (P)$ và $b\perp a$ thì $b\parallel (P)$ hoặc $b\subset (P)$}
		{\True Nếu $a\parallel (P)$  và $b\perp a$ thì $b\perp (P)$}
	\loigiai{ 
		Mệnh đề sai là \lq\lq Nếu $a\parallel (P)$ và $b\perp a$ thì $b\perp (P)$\rq\rq, vì đường  thẳng $b$ có thể nằm trong mặt phẳng $(P)$.
	 }
\end{ex}

% nguồn: Ok_Toan11_TracNghiem_Lan2
\begin{ex}%[1H8N2-1]%[ma:1cbd2c63]
	Khẳng định nào  sau đây là đúng?
	 \choice
		{Hai đường thẳng cùng vuông góc  với đường thẳng thứ ba thì chúng vuông góc với nhau}
		{Hai đường thẳng cùng vuông góc với một mặt phẳng  thì chúng song song với nhau}
		{\True Hai đường thẳng phân biệt  cùng vuông góc với một mặt phẳng thì chúng song song với nhau}
		{Hai đường thẳng phân biệt cùng vuông góc với một  đường thẳng thì chúng song song với nhau}
	\loigiai{ 
		Khẳng định đúng là \lq\lq Hai đường thẳng phân biệt cùng vuông góc với một mặt phẳng  thì chúng song song với nhau\rq\rq.
	 }
\end{ex}

% nguồn: Ok_Toan11_TracNghiem_Lan2
\begin{ex}%[1H8N2-1]%[ma:bb437514]
	Cho hai đường thẳng $a$, $b$ phân  biệt và mặt phẳng $(P)$. Mệnh đề nào sau đây sai?
	\choice 
		{Nếu $(P)\parallel (Q)$  và $b\perp (P)$ thì $b\perp (Q)$}
		{\True Nếu $a\parallel (P)$ và $b\perp a$ thì $b\perp (P)$} 
		{Nếu $a\parallel  (P)$ và $b\perp (P)$ thì $b\perp a$}
		{Nếu $a\perp (P)$ và $b\perp  (P)$ thì $a\parallel b$}
	\loigiai{ 
		Mệnh đề sai là \lq\lq Nếu $a\parallel (P)$ và $b\perp a$ thì $b\perp (P)$\rq\rq. 
	} 
\end{ex}

% nguồn: Ok_Toan11_TracNghiem_Lan2
\begin{ex}%[1H8N2-1]%[ma:1e4504e7]
	Trong không gian cho điểm $O$ và  đường thẳng $d$. Qua điểm $O$ có bao nhiêu mặt phẳng vuông góc với đường thẳng $d$?
	 \choice
		{Ba} 
		{Hai} 
		 {\True Một}
		 {Vô số}
	\loigiai{ 
		Theo tính chất ta có: Qua điểm $O$ có một mặt phẳng vuông góc với đường thẳng  $d$.
	 }
\end{ex}

% nguồn: Ok_Toan11_TracNghiem_Lan2
\begin{ex}%[1H8N2-1]%[ma:0871f796]
	Qua điểm $O$ cho trước, có bao nhiêu mặt phẳng vuông góc với  đường thẳng $\Delta$ cho trước?
	\choice 
		 {\True $1$}
		{Vô  số}
		{$3$} 
		{$2$} 
	\loigiai{ 
		Theo tính chất 1: Có duy nhất một mặt phẳng đi qua một  điểm cho trước và vuông góc với một đường thẳng cho trước.
	 }
\end{ex}

% nguồn: Ok_Toan11_TracNghiem_Lan2
\begin{ex}%[1H8N2-1]%[ma:e4e7a8ee]
	Mệnh đề nào sau đây là  đúng?
	 \choice
		{Một đường thẳng vuông góc với một trong hai đường thẳng vuông góc thì song song  với đường thẳng còn lại}
		{Hai đường thẳng  cùng vuông góc với một đường thẳng thì song song với nhau}
		{Hai đường thẳng cùng vuông góc với một đường thẳng thì  vuông góc với nhau}
		{\True Một đường thẳng vuông  góc với một trong hai đường thẳng song song thì vuông góc với đường thẳng kia}
	 \loigiai{
		Mệnh đề đúng là \lq\lq Một đường thẳng vuông góc với một trong hai đường thẳng song song  thì vuông góc với đường thẳng kia\rq\rq.
	} 
\end{ex}

% nguồn: Ok_TracNghiemLop11_Lan1
\begin{ex}%[1H8N2-1]%[ma:3e18df29]
	Khẳng định nào sau  đây \textbf{sai}?
	 \choice
		{Nếu đường thẳng $d$ vuông góc với mặt phẳng $\left( \alpha \right)$ thì $d$ vuông góc với hai đường thẳng trong mặt  phẳng $\left( \alpha \right)$}
		{\True Nếu đường thẳng $d$ vuông góc với hai đường thẳng nằm trong mặt phẳng $\left( \alpha \right)$ thì $d$  vuông góc với mặt phẳng $\left( \alpha \right)$}
		{Nếu đường thẳng $d$ vuông góc với hai đường thẳng cắt nhau nằm trong mặt phẳng $\left( \alpha \right)$ thì $d$ vuông góc với bất  kỳ đường thẳng nào nằm trong mặt phẳng $\left( \alpha \right)$}
		{Nếu $d\perp \left( \alpha \right)$ và đường  thẳng $a\,\parallel \left( \alpha \right)$ thì $d\perp a$}
	 \loigiai{
	Khẳng định $B$ sai vì: đường thẳng $d$ vuông góc với hai đường thẳng nằm trong mặt  phẳng $\left( \alpha \right)$ mà hai đường thẳng đó song song thì $d$ không vuông góc với mặt phẳng $\left( \alpha \right)$.}
\end{ex}

% nguồn: Ok_TracNghiemLop11_Lan1
\begin{ex}%[1H8N2-1]%[ma:527e10b3]
	Cho hai đường thẳng phân biệt $a, b$ và mặt phẳng $(P)$. Chọn  khẳng định đúng?
	 \choice
		{Nếu  $a \parallel (P)$ và $b \perp a$ thì $b \perp(P)$}
		{\True Nếu $a \parallel (P)$ và $b  \perp(P)$ thì $b \perp a$}
		{Nếu $a \perp(P)$ và $b \perp a$ thì  $b \parallel (P)$}
		{Nếu $a \parallel (P)$ và $b \parallel (P)$ thì $b  \parallel a$}
	\loigiai{ 
		"Nếu $a \parallel (P)$ và $b \perp(P)$ thì $b \perp a$" là khẳng  định đúng.
	 }
\end{ex}

% nguồn: Ok_TracNghiemLop11_Lan1
\begin{ex}%[1H8N2-1]%[ma:ec5a70fe]
	Trong không gian cho ba đường thẳng phân biệt $a$, $b$, $c$. Khẳng  định nào sau đây đúng?
	 \choice
		{ Nếu  $a$ và $b$ cùng nằm trong một mặt phẳng và cùng vuông góc với $c$ thì $a\parallel b$}
		 {\True Nếu $a\parallel b$ và $c\perp a$ thì $c\perp b$}
		{Nếu $a$, $b$, $c$ cùng nằm trong mặt phẳng $(P)$ và góc giữa  $a$ và $c$ bằng góc giữa $b$ và $c$ thì $a\parallel b$}
		{Nếu $a$ và $b$ cùng nằm trong mặt phẳng $(P)$  và $(P)\parallel c$ thì góc giữa $a$ và $c$ bằng góc giữa $b$ và $c$}
	 \loigiai{
		Khẳng định đúng là ``Nếu  $a\parallel b$ và $c\perp a$ thì $c\perp b$''.
	 }
\end{ex}

% nguồn: Ok_TracNghiemLop11_Lan1
\begin{ex}%[1H8N2-1]%[ma:a6de1321]
	Trong không gian cho điểm  $A$ và đường thẳng $d$. Có bao nhiêu đường thẳng qua $A$ và vuông góc với đường thẳng $d$?
	\choice 
		 {$0$}
		 {$1$}
		{\True  Vô số}
		{$2$} 
	\loigiai{ 
	Vì qua điểm $A$ có thể dựng vô số đường thẳng vuông góc với $d$, các đường  thẳng đó cùng nằm trong một mặt phẳng vuông góc với $d$.}
\end{ex}

% nguồn: Ok_TracNghiemLop11_Lan1
\begin{ex}%[1H8N2-1]%[ma:8cb130db]
	Trong không gian, cho  hai đường thẳng $a$, $b$ , và mặt phẳng $(P)$. Trong các khẳng định sau,
	 khẳng định nào là \textbf{đúng}?
	 \choice
		 {\True $\heva{&a\parallel b\\&a\perp(P)}\Rightarrow b\perp(P)$}
		 {$\heva{&a\perp b\\&b\parallel(P)}\Rightarrow a\perp(P)$}
		 {$\heva{&a\perp b\\&b\perp(P)}\Rightarrow a\perp(P)$}
		{$\heva{&a\perp  b\\&b\perp(P)}\Rightarrow a\parallel(P)$}
	 \loigiai{
	 }
\end{ex}

% nguồn: Ok_TracNghiemLop11_Lan1
\begin{ex}%[1H8N2-1]%[ma:b4be979b]
	Chỉ ra mệnh đề sai trong các mệnh đề  sau:
	\choice 
		{Qua một điểm  $O$ cho trước có một và chỉ một đường thẳng vuông góc với một mặt phẳng cho trước}
		{Qua một điểm $O$ cho trước có một mặt phẳng duy nhất vuông góc với một đường thẳng $\Delta$ cho  trước}
		{Hai đường thẳng chéo nhau và vuông góc với nhau. Khi đó có một và chỉ một mặt phẳng chứa đường thẳng này và vuông góc với  đường thẳng kia}
		{\True Qua một điểm $O$ cho trước có một và chỉ một đường thẳng vuông góc với một đường thẳng  cho trước}
	 \loigiai{
	Qua một điểm $O$ cho trước có vô số đường thẳng vuông góc với một  đường thẳng cho trước, chúng cùng nằm trên mặt phẳng qua $O$ và vuông góc với đường thẳng cho trước.}
\end{ex}

% nguồn: Ok_TracNghiemLop11_Lan1
\begin{ex}%[1H8N2-1]%[ma:de67d94a]
	 Trong các mệnh đề sau, mệnh đề nào \textbf{sai}?
	\choice 
		{Cho đường thẳng $a$ vuông  góc với đường thẳng $b$ và $b$ nằm trong mặt phẳng $(P)$; mọi mặt phẳng $(Q)$ chứa $a$ và vuông góc với $b$ thì $(P)$ vuông góc với $(Q)$}
		{\True Nếu đường thẳng $a$ vuông góc với đường thẳng $b$  và mặt phẳng $(P)$ chứa $a$, mặt phẳng $(Q)$ chứa $b$ thì $(P)$ vuông góc với $(Q)$}
		{Cho đường thẳng $a$ vuông góc với mặt phẳng $(P)$, mọi mặt phẳng  $(Q)$ chứa $a$ thì $(P)$ vuông góc với $(Q)$}
		{Qua một điểm có duy nhất một mặt  phẳng vuông góc với một đường thẳng cho trước}
	 \loigiai{
	\immini{ Ta thấy rằng hai mặt phẳng trong hình lần lượt chứa hai đường thẳng  vuông góc $a$ và $b$ nhưng chúng không vuông góc với nhau. Góc $\varphi$ có thể nhận một giá trị khác $90^{\circ}$.
		} 
		 {
		\begin{tikzpicture}[scale=0.7,font=\footnotesize,line join=round,line cap=round,>=stealth]
				\path 
					(0,0) coordinate (A) 
					(-2,-2) coordinate (B) 
					(3,-2)  coordinate (C)
					($(A)+(C)-(B)$)  coordinate (D)
					 ($(B)!0.6!(C)$) coordinate (E)
					($(A)!0.6!(D)$) coordinate  (F)
					($(B)!0.3!(E)$)  coordinate (M)
					($(M)+(0,4)$) coordinate  (N)
					($(F)+(N)-(E)$)  coordinate (K)
					($(E)!0.5!(F)$) coordinate  (L)
					($(B)+(L)-(E)$)  coordinate (H)
					 ($(N)+(L)-(E)$) coordinate (G)
					($(G)!0.3!(L)$) coordinate  (I)
					($(H)!0.6!(L)$) coordinate  (J);
				\draw (A)--(B)--(C)--(D)--(F) 
					 (N)--(E)--(F) --(K)--(N)
					(I)--(L); 
				\draw[dashed]  (A)--(F)
					(J)--(L); 
				 \draw pic[draw,angle radius=0.3cm] {angle = I--L--J};
				\draw pic[draw,angle radius=0.2cm]  {right angle = I--L--F};
				\draw pic[draw,angle radius=0.2cm] {right angle =  J--L--E};
			\end{tikzpicture}
		 }
	} 
\end{ex}

% nguồn: Ok_TracNghiemLop11_Lan1
\begin{ex}%[1H8N2-1]%[ma:07dce318]
	Cho điểm $M$ và mặt phẳng $(\alpha)$ có bao nhiêu đường thẳng đi qua điểm $M$ và vuông  góc với mặt phẳng $(\alpha)$?
	\choice 
		 {$2$}
		{Vô  số}
		 {$0$}
		{\True $1$} 
	\loigiai{ 
		Theo tính chất qua một điểm cho trước chỉ có duy  nhất một đường thẳng vuông góc với một mặt phẳng cho trước.
	 }
\end{ex}

% nguồn: Ok_TracNghiemLop11_Lan1
\begin{ex}%[1H8N2-1]%[ma:7e8d017d]
	Đường thẳng $d$  vuông góc với mặt phẳng $(\alpha)$ nếu đường thẳng $d$ vuông góc với
	\choice 
		{hai đường thẳng $a$ và $b$ cùng nằm trong mặt  phẳng $(\alpha)$}
		{một  đường thẳng $a$ nằm trong mặt phẳng $(\alpha)$}
		 {hai đường thẳng cắt nhau $a$ và $b$}
		{\True mọi đường  thẳng $a$ nằm trong mặt phẳng $(\alpha)$}
	 \loigiai{
		Đường thẳng $d$ vuông góc với mặt phẳng $(\alpha)$ nếu đường thẳng $d$ vuông góc với mọi đường thẳng $a$ nằm trong mặt phẳng  $(\alpha)$.
	} 
\end{ex}

% nguồn: Ok_TracNghiemLop11_Lan1
\begin{ex}%[1H8N2-1]%[ma:c82a78e2]
	Cho điểm $M$ và đường thẳng $a$ có bao nhiêu mặt  phẳng đi qua điểm $M$ và vuông góc với đường thẳng $a$?
	\choice 
		{$2$} 
		{Vô số} 
		 {$0$}
		{\True $1$} 
	\loigiai{ 
		Theo tính chất qua một điểm cho trước chỉ có  duy nhất một mặt phẳng vuông góc với đường thẳng cho trước.
	 }
\end{ex}

% nguồn: Ok_TracNghiemLop11_Lan1
\begin{ex}%[1H8N2-1]%[ma:9cb0b422]
	Trong không gian, có bao nhiêu mặt phẳng qua điểm $O$ cho trước và vuông góc với đường thẳng $\Delta$ cho  trước?
	\choice 
		 {$3$}
		{\True  $1$}
		 {$2$}
		{Vô  số}
	\loigiai{ 
		Chỉ có duy nhất một mặt phẳng qua điểm $O$ cho trước  và vuông góc với đường thẳng $\Delta$ cho trước
	} 
\end{ex}

% nguồn: Ok_TracNghiemLop11_Lan1
\begin{ex}%[1H8N2-1]%[ma:af1b8aa8]
	Trong các mệnh đề sau đây,  mệnh đề nào đúng?
	\choice 
		{\True  Hai đường thẳng phân biệt cùng vuông góc với đường thẳng thứ ba thì có thể vuông góc với nhau}
		{Hai đường thẳng phân biệt cùng vuông góc với đường thẳng thứ ba  thì song song với nhau}
		{Hai đường thẳng phân  biệt cùng vuông góc với đường thẳng thứ ba thì vuông góc với nhau}
		{Hai đường thẳng phân biệt cùng vuông góc với đường thẳng thứ  ba thì hoặc song song với nhau hoặc vuông góc với nhau}
	\loigiai{ 
		Hai đường thẳng phân biệt cùng vuông góc với đường thẳng  thứ ba thì có thể vuông góc với nhau.
	} 
\end{ex}

% nguồn: Ok_TracNghiemLop11_Lan1
\begin{ex}%[1H8N2-1]%[ma:047e907f]
	Cho tam giác $ABC$. Có bao nhiêu  mặt phẳng đi qua $A$ và vuông góc với đường thẳng $AB$?
	 \choice
		 {\True $1$}
		 {$2$}
		 {$3$}
		{Vô  số}
	 \loigiai{
		Có $1$ mặt phẳng đi qua $A$ và vuông góc với  đường thẳng $AB$.
	} 
\end{ex}

% nguồn: Ok_TracNghiemLop11_Lan1
\begin{ex}%[1H8N2-1]%[ma:fd943c34]
	Trong không gian, cho $3$ đường thẳng $a$, $b$, $c$ phân biệt và mặt phẳng $(P)$.  Mệnh đề nào sau đây đúng?
	 \choice
		{Nếu $a  \perp c$ và $(P) \perp c$ thì $a \parallel (P)$}
		{Nếu $a \perp c$ và $b  \perp c$ thì $a \parallel b$}
		 {Nếu $a \perp b$ và $b \perp c$ thì $a \perp c$}
		{\True Nếu $a \perp b$ thì $a$  và $b$ cắt nhau hoặc chéo nhau}
	 \loigiai{
		Có bốn vị trí tương đối của hai đường thẳng mà nếu hai đường thẳng  trùng nhau hoặc song song thì chúng không vuông góc với nhau do đó nếu $a \perp b$ thì $a$ và $b$ cắt nhau hoặc chéo nhau.
	 }
\end{ex}

% nguồn: Ok_TracNghiemLop11_Lan1
\begin{ex}%[1H8N2-1]%[ma:1309065e]
	Trong  không gian cho đường thẳng $\Delta$ và điểm $O$. Qua $O$ có mấy đường thẳng vuông góc với $\Delta$ cho trước?
	 \choice
		 {$1$}
		{$2$} 
		{$3$} 
		 {\True Vô số}
	\loigiai{Qua $O$ có vô số đường thẳng vuông góc với  $\Delta$ cho trước, tất cả các đường thẳng này cùng nằm trong mặt phẳng qua $O$ và vuông góc với đường thẳng $\Delta$ cho trước.}
\end{ex}

% nguồn: Ok_TracNghiemLop11_Lan1
\begin{ex}%[1H8N2-1]%[ma:e22412b6]
	Mệnh  đề nào đúng trong các mệnh đề sau đây?
	\choice 
		{Góc giữa đường thẳng $a$ và mặt phẳng $(P)$ bằng góc giữa đường thẳng $a$ và mặt phẳng $(Q)$ thì mặt phẳng  $(P)$ song song hoặc trùng với mặt phẳng $(Q)$}
		{Góc giữa đường thẳng $a$ và mặt phẳng $(P)$ bằng góc giữa đường thẳng $b$ và  mặt phẳng $(P)$ thì đường thẳng $a$ song song với đường thẳng $b$}
		{Góc giữa đường  thẳng $a$ và mặt phẳng $(P)$ bằng góc giữa đường thẳng $b$ và mặt phẳng $(P)$ thì đường thẳng $a$ song song hoặc trùng với đường thẳng $b$}
		{\True Góc giữa đường thẳng và mặt phẳng bằng góc giữa đường thẳng đó và hình chiếu của nó trên mặt phẳng đã cho. (Biết đường thẳng không vuông góc với  mặt phẳng)}
	 \loigiai{
		Khi đường thẳng không vuông góc với mặt phẳng thì góc giữa đường thẳng và mặt  phẳng bằng góc giữa đường thẳng đó và hình chiếu của nó trên mặt phẳng đã cho.
	 }
\end{ex}

% nguồn: Ok_TracNghiemLop11_Lan1
\begin{ex}%[1H8N2-1]%[ma:ff32a06d]
	Trong không gian cho trước điểm $M$ và đường thẳng $\Delta $. Các đường thẳng đi qua  $M$ và vuông góc với $\Delta $ thì:
	 \choice
		{Vuông góc với  nhau}
		{Song  song với nhau}
		{Cùng vuông góc với một mặt phẳng} 
		{\True Cùng  thuộc một mặt phẳng}
	 \loigiai{
	Suy ra từ tính chất $1$ theo SGK hình học $11$ trang $100$.} 
\end{ex}

% nguồn: Ok_TracNghiemLop11_Lan1
\begin{ex}%[1H8H2-1]%[ma:9c0dfa87]
	Cho hình chóp $S.ABCD$ có đáy $ABCD$ là hình thoi tâm $O$, $SA\perp (ABCD)$. Khẳng định nào sau đây  là \textbf{sai}?
	\choice 
		{$SA  \perp BD $}
		{$SO \perp BD $} 
		{\True  $AD \perp SC $}
		{$SC  \perp BD $}
	\loigiai{ 
	\immini{\begin{itemize}
		\item Vì $SA\perp (ABCD)$ nên $SA\perp BD$, do đó khẳng định này  đúng.
		\item Vì tam giác $SBD$ cân tại $S$ và $O$ là trung điểm $BD$ nên $SO\perp BD$, do đó khẳng định  này đúng.
		\item  Vì $BD\perp AB, BD\perp SA$ nên $BD\perp (SAC)$, suy ra $BD\perp SC$, do đó khẳng định này đúng.
		\item Khẳng định sai là $AD \perp  SC$.
	\end{itemize}}{\begin{tikzpicture}[scale=.75,font=\footnotesize,line join=round,line cap=round,>=stealth]
		\path (0:0) coordinate (A)++(0:4) coordinate (D)++(-120:2) coordinate (C) ++(180:4)  coordinate (B)
			 (A)++(90:3) coordinate (S)
			($(A)!.5!(C)$)  coordinate (O);
		\draw[dashed]  (S)--(A)--(B)--(D)--(A)--(C);
		\draw  (S)--(B)--(C)--(S)--(D)--(C);
		\foreach \x/\g in  {S/90,A/150,B/-90,C/-90,D/0,O/90} \fill (\x) circle (1pt)+(\g:.25) node {$\x$};
	\end{tikzpicture}
	 }
} 
\end{ex}

% nguồn: Ok_TracNghiemLop11_Lan1
\begin{ex}%[1H8H2-1]%[ma:dd969808]
	Trong  các mệnh đề sau, mệnh đề nào \textbf{sai}?
	\choice 
		{Hai  đường thẳng vuông góc nếu góc giữa hai đường thẳng bằng $90^{\circ}$}
		{Hai đường thẳng phân biệt cùng song song với một  đường thẳng thì song song với nhau}
		{Một đường thẳng vuông  góc với hai cạnh của tam giác thì sẽ vuông góc với cạnh thứ ba của tam giác đó}
		{\True  Hai đường thẳng phân biệt cùng vuông góc với một mặt phẳng thì vuông góc với nhau}
	\loigiai{ 
		Mệnh đề sai  là \lq\lq Hai đường thẳng phân biệt cùng vuông góc với một mặt phẳng thì vuông góc với nhau\rq\rq.
	} 
\end{ex}

% nguồn: Ok_TracNghiemLop11_Lan1
\begin{ex}%[1H8H2-1]%[ma:10dde67b]
	Cho hai đường thẳng phân biệt $a$, $b$ và mặt phẳng $(P)$, trong đó $a\perp(P)$. Chọn mệnh đề  sai.
	 \choice
		{Nếu $b\parallel a$  thì $b\parallel(P)$}
		{\True Nếu $b\parallel  a$ thì $b\perp(P)$}
		{Nếu $b\perp(P)$  thì $b\parallel a$}
		 {Nếu $b\parallel(P)$ thì $b\perp a$}
	\loigiai{ 
		Nếu $a\perp(P)$ và $b\parallel a$  thì $b\perp(P)$.
	 }
\end{ex}

% nguồn: Ok_TracNghiemLop11_Lan1
\begin{ex}%[1H8H2-1]%[ma:8c9346b5]
	Trong không gian, mệnh đề nào sau  đây đúng?
	\choice 
		{Nếu hai đường  thẳng vuông góc với nhau thì hai đường thẳng đó cắt nhau}
		 {Nếu hai đường thẳng vuông góc với đường thẳng thứ ba thì hai đường thẳng đó song song với nhau}
		{Nếu hai đường thẳng vùng vuông góc với đường thẳng thứ  ba thì hai đường thẳng đó vuông góc với nhau}
		{\True Cho  hai đường thẳng song song với nhau. Nếu một đường thẳng vuông góc với đường thẳng này thì cũng vuông góc với đường thẳng kia}
	\loigiai{ 
		Theo  tính chất liên hệ giữa quan hệ song song và quan hệ vuông góc trong mặt phẳng.
	} 
\end{ex}

% nguồn: Ok_TracNghiemLop11_Lan1
\begin{ex}%[1H8H2-1]%[ma:d8493a68]
	 Mệnh đề nào sau đây \textbf{sai}?
	\choice 
		{Hai đường thẳng phân biệt cùng vuông góc với một mặt phẳng thì song  song}
		{Hai mặt phẳng phân  biệt cùng vuông góc với một đường thẳng thì song song}
		{\True Cho đường thẳng $\Delta$ song song với  mặt phẳng $(\alpha)$}
		{Một đường thẳng và một mặt phẳng (không chứa đường  thẳng đã cho) cùng vuông góc với một đường thẳng thì song song nhau}
	\loigiai{ 
		C sai vì đường thẳng vuông góc với $\Delta$ có thể song song với $(\alpha)$ nếu nó nằm trong cùng một mặt phẳng với  đường thẳng $\Delta$.\\
	} 
\end{ex}

% nguồn: Ok_TracNghiemLop11_Lan1
\begin{ex}%[1H8H2-1]%[ma:858954ef]
	Cho tam giác $A B C$ có bao nhiêu mặt phẳng đi qua điểm $A$ và vuông góc với  đường thẳng $A B$?
	 \choice
		{$2$} 
		 {Vô số}
		 {$0$}
		{\True  $1$}
	\loigiai{ 
		Theo tính chất qua  một điểm cho trước chỉ có duy nhất một mặt phẳng vuông góc với đường thẳng cho trước.
	 }
\end{ex}

% nguồn: Ok_TracNghiemLop11_Lan1
\begin{ex}%[1H8H2-1]%[ma:082ca765]
	Cho hình bình hành $A B C D$ tâm $O$ có bao nhiêu đường thẳng đi qua điểm $O$ và vuông góc với mặt phẳng $(A B C  D)$?
	 \choice
		 {$2$}
		{\True  $1$}
		 {Vô số}
		{$0$} 
	 \loigiai{
		Theo tính chất qua một điểm cho trước  chỉ có duy nhất một đường thẳng vuông góc với một mặt phẳng cho trước.
	} 
\end{ex}

% nguồn: Ok_TracNghiemLop11_Lan1
\begin{ex}%[1H8H2-1]%[ma:62060d04]
	\immini{Cho đường thẳng $a$ nằm trong mặt phẳng $(P)$, $b$ là đường thẳng không nằm trong $(P)$ và không vuông góc với $(P)$. Gọi $b'$ là hình chiếu vuông góc của $b$ trên $(P)$ (tham khảo hình  bên).
	 \choice
		{$a$ vuông góc với $b'$ khỉ và chi khi $a$ vuông góc  với $(P)$}
		{$a$ vuông góc với $b$ khi và chỉ  khi $a$ vuông góc với $(P)$}
		{$a$ vuông góc  với $b'$ khi và chỉ khi $b$ vuông góc với $b'$}
		{\True $a$ vuông góc với $b$ khi và chỉ khi $a$  vuông góc với $b'$}}
		{\begin{tikzpicture}[scale=.6,font=\footnotesize,line join=round,line cap=round,>=stealth]
		\path 
			(0,0) coordinate (A) 
			(7,0) coordinate  (B)
			(60:3.5) coordinate  (C)
			 ($(B)+(C)-(A)$) coordinate (D)
			(3.5,4.5) coordinate (M) 
			(3.5,1.5) coordinate  (N)
			 (6.5,1.8) coordinate (K)
			(6.5,4) coordinate (H) 
			($(H)!1.3!(M)$) coordinate (A1) 
			($(M)!1.3!(H)$) coordinate (A2) 
			($(K)!1.3!(N)$) coordinate  (A3)
			($(N)!1.3!(K)$) coordinate  (A4)
			(intersection of M--N and C--D)  coordinate (O1)
			(intersection of H--K and C--D)  coordinate (O2);
		\draw  (C)--(A)--(B)--(D)
			(C)--(O1) 
			(D)--(O2) 
			(A3)--(A4) 
			 (M)--(N)
			 (K)--(H)
			pic[draw,angle  radius=6mm]{angle=B--A--C};
		\draw (4.5,0.5)--(5.5,2.5)  node[left]{$a$};
		 \draw (A1)--(A2);
		 \draw[dashed] (O1)--(O2);
		\path  (A)+(30:4.5mm) node{$P$}
			(M)--(A1)  node[midway,above]{$b$}
			(N)--(A3) node[midway,above]{$b'$}; 
	\end{tikzpicture}
		} 
	\loigiai{ 
		Theo định lý ba đường vuông góc ta có $a$ vuông góc với $b$  khi và chỉ khi $a$ vuông góc với $b'$.
	 }
\end{ex}

% nguồn: Ok_TracNghiemLop11_Lan1
\begin{ex}%[1H8H2-1]%[ma:507c751a]
	Cho hai đường thẳng phân biệt $a, b$ và mặt phẳng $(P)$, trong đó $a\perp(P)$. Mệnh  đề nào sau đây là \textbf{sai}?
	\choice 
		{Nếu $b \perp(P)$ thì  $b \parallel a$}
		{Nếu  $b \parallel (P)$ thì $b \perp a$}
		 {Nếu $b \parallel a$ thì $b\perp(P)$}
		{\True Nếu  $b\perp$ thì $b \parallel(P)$}
	\loigiai{ 
		D sai vì đường thẳng $b$  có thể nằm trong mặt phẳng $(P)$.\\
	 }
\end{ex}

% nguồn: Ok_TracNghiemLop11_Lan1
\begin{ex}%[1H8H2-1]%[ma:9659c9fd]
	Trong không gian, mệnh đề nào sau  đây đúng?
	 \choice
		{\True Nếu $a\parallel (P)$ và $b\perp (P)$  thì $b\perp a$}
		{Nếu $a\parallel (P)$ và $b\parallel a$ thì $b\parallel  (P)$}
		{Một đường thẳng  vuông góc với hai đường thẳng phân biệt trong mặt phẳng $(P)$ thì nó vuông góc với mặt phẳng $(P)$}
		{Nếu $a\parallel (P)$ và $b\perp a$ thì $b\perp (P)$} 
	 \loigiai{
	Theo tính chất liên hệ giữa  quan hệ song song và quan hệ vuông góc.}
\end{ex}

% nguồn: Ok_TracNghiemLop11_Lan1
\begin{ex}%[1H8V2-1]%[ma:2d646fff]
	 Trong không gian tập hợp các điểm $M$ cách đều hai điểm cố định $A$ và $B$ (phân biệt) là
	\choice{\True Mặt phẳng trung trực của đoạn thẳng  $AB$}
		{Đường  trung trực của đoạn thẳng $AB$}
		{Mặt  phẳng vuông góc với $AB$ tại $A$}
		{Đường thẳng qua  $A$ và vuông góc với $AB$}
	\loigiai{ 
		Tập hợp các điểm $M$  cách đều hai điểm cố định $A$ và $B$ là mặt phẳng trung trực của đoạn thẳng $AB$.
	 }
\end{ex}


%%%%%%%%%%%%%%%% ĐÚNG/SAI %%%%%%%%%%%%%%%%

% nguồn: Ok_dungsai_lop11_Lan1
\begin{ex}%[1H8N2-1]%[ma:0b88a237]
 Xét tính Đúng/ Sai các khẳng  định sau
 \choiceTF 
  {Có vô số đường thẳng đi qua một điểm cho trước  và vuông góc với mặt phẳng cho trước}
  {\True Đường thẳng vuông góc với một mặt phẳng thì vuông góc với mọi đường thẳng nằm trong mặt  phẳng đó}
  {Nếu một đường thẳng vuông góc với hai đường thẳng cùng nằm  trong một mặt phẳng thì nó vuông góc với mặt phẳng ấy}
  {\True Có duy nhất một mặt phẳng  đi qua một điểm cho trước và vuông góc với đường thẳng cho trước}
 \loigiai{ 
  \begin{itemchoice}
   \itemch  Sai
    \itemch Đúng
    \itemch Sai
   \itemch  Đúng
  \end{itemchoice}
  }
\end{ex}

% nguồn: Ok_dungsai_lop11_Lan1
\begin{ex}%[1H8H2-1]%[ma:88c444c1]
 Cho hai  đường thẳng phân biệt $a$, $b$ và mặt phẳng $(P)$, trong đó $a\perp(P)$.
 \choiceTF 
  {Nếu $b\parallel a$ thì $b\parallel(P)$} 
  {\True Nếu $b\parallel  a$ thì $b\perp(P)$}
  {Nếu $b\perp(P)$ thì $b\parallel  a$}
  {Nếu $b\parallel(P)$ thì $b\perp  a$}
 \loigiai{ 
   Nếu $a\perp(P)$ và $b\parallel a$ thì $b\perp(P)$.
 } 
\end{ex}
